% Drop-in subsection for the experiments/results section.
% Numerical source of record: analyses/probe_validity/paper/paper_results.csv
% Required packages already used by the manuscript: booktabs, graphicx.

\subsection{Probe validity through capability decomposition}
\label{sec:probe-validity}

\paragraph{Motivation and diagnostic setup.}
A high score on a learned probe is not by itself evidence for the intended
capability.  In one-step feature prediction (C1), retrieval can benefit from
representation continuity or copying the last observation even when the
predictor makes little use of the camera action.  Likewise, low hidden-object
localization error (B2) need not imply an ordered spatial-memory update: it may
reflect the appearance query, an unordered collection of historical features,
or dataset regularities.  We therefore decompose each probe with
inference-time interventions while keeping the trained predictor and target
fixed.  These tests measure information use by a fixed predictor; they are not
retrained ablations and do not establish necessity, planning ability, or a
complete dynamics model.

For C1, we compare the prediction under the correct 9D relative camera action
with one under a wrong action matched for total pose displacement.  The primary
pool contains eight physical candidates
(\textsc{Matched-Total-Pose-Displacement}, candidate seed 42), not the
same-scene pool, which contains only four unique candidates.  We report
\begin{equation}
  \Delta_{\mathrm{action}}
  = R@1_{\mathrm{correct}}-R@1_{\mathrm{matched\ wrong}}.
\end{equation}
For B2, an object-reference query is pooled using a ground-truth object mask
from a historical visible frame.  The spatial readout receives that appearance
query and the historical feature tokens, but neither the mask nor camera pose;
thus the task is not mask-free at query construction.  Perturbation effects are
paired on the samples for which each condition is defined:
\begin{equation}
  \Delta_{\mathrm{perturb}}=E_{\mathrm{perturb}}-E_{\mathrm{full}},
\end{equation}
where $E$ is spherical angular error in degrees and positive values indicate
degradation.

\paragraph{C1: predictability is distinct from action sensitivity.}
Table~\ref{tab:probe-validity-c1} separates global predictability from action
sensitivity.  At p12, Wan has the highest global correct-action R@1 (15.22\%),
whereas V-JEPA2 has the largest hard-pool action gain (18.84 points).  The
resulting model order is therefore not the order induced by global retrieval.
Moreover, copying is stronger than correct-action global retrieval for
V-JEPA2 and VideoMAE at both prefixes, yet their hard-pool gains remain
positive.  Copy-based predictability and action use can consequently coexist.

The p12 scene-clustered intervals exclude zero for V-JEPA2, VideoMAE, and
CogVideoX, but include zero for Wan; the verified p16 repair evaluation
shows the same qualitative pattern.  V-JEPA2 has the largest point estimate at
both prefixes, but we do not interpret separate within-model intervals as a
formal test of between-model differences.  The p12 and p16 samples are also
not strictly target-paired, so their agreement is cross-prefix-setting
robustness rather than a causal effect of adding history.

\begin{table*}[t]
  \centering
  \caption{\textbf{C1 predictability and action sensitivity.} Global columns
  report correct-action/copy R@1 in percent.  $\Delta_{\mathrm{action}}$ is the
  R@1 difference in percentage points for the physical
  \textsc{Matched-Total-Pose-Displacement} pool with $K=8$ and candidate seed
  42.  Each prefix has $n=276$ samples from 69 scenes and one probe-training
  seed (42).  Brackets are 95\% scene-clustered bootstrap intervals and do not
  include probe-training seed uncertainty.  The p12 and p16 target sets are not
  strictly paired.}
  \label{tab:probe-validity-c1}
  \small
  \setlength{\tabcolsep}{4.5pt}
  \begin{tabular}{@{}lcrrrr@{}}
    \toprule
    & & \multicolumn{2}{c}{p12} & \multicolumn{2}{c}{p16} \\
    \cmidrule(lr){3-4}\cmidrule(l){5-6}
    Model & Layer & Global correct/copy & $\Delta_{\mathrm{action}}$ [95\% CI] &
      Global correct/copy & $\Delta_{\mathrm{action}}$ [95\% CI] \\
    \midrule
    V-JEPA2   & 23 & 13.04/19.93 & $+18.84\ [13.04,25.36]$ & 13.41/23.91 & $+18.84\ [13.41,23.91]$ \\
    VideoMAE  & 23 &  9.42/17.39 & $+10.51\ [ 6.16,15.22]$ &  5.07/ 9.06 & $+10.51\ [ 5.43,15.58]$ \\
    Wan       & 20 & 15.22/11.96 & $ +2.90\ [-1.45, 7.25]$ & 10.51/15.22 & $ +1.81\ [-2.17, 5.80]$ \\
    CogVideoX & 20 &  3.26/ 2.54 & $ +4.35\ [ 0.35, 7.97]$ &  2.90/ 1.45 & $ +6.52\ [ 2.54,10.87]$ \\
    \bottomrule
  \end{tabular}
\end{table*}

\paragraph{B2: dependence does not establish ordered state updating.}
Table~\ref{tab:probe-validity-b2} shows that query and history dependence are
model-specific.  Masking the query or swapping it within the same scene
reliably degrades V-JEPA2 and VideoMAE.  Wan instead shows clearer dependence
on historical content: retaining only the current frame increases error by
$10.01^{\circ}$, and removing frames in which the target was visible increases
it by $6.59^{\circ}$.  CogVideoX degrades under query swapping, but not under
query masking or the current-frame intervention.  Because these are
fixed-predictor inputs, degradation can reflect distribution shift and is not
equivalent to a query-only or current-frame-only probe trained from scratch.

Content-only shuffling and reversal leave the query, ground truth, final
reference frame, and frame identifiers fixed.  Their clustered intervals
include zero for every model.  Thus, \emph{under the evaluated protocol, we do
not detect a reliable content-slot order effect}.  This does not show that
order is unimportant, that a backbone lacks temporal information, or that B2
reduces to unordered retrieval.  In particular, temporal information may
already be encoded within each feature.  Moving content together with its
original frame identifier produces exactly zero change and is retained only as
a structural sanity check, not as the primary scientific test.

\begin{table*}[t]
  \centering
  \caption{\textbf{B2 fixed-predictor diagnostics at p12.} FULL is mean
  spherical angular error in degrees ($n=165$, 65 scenes).  Other entries are
  paired $\Delta_{\mathrm{perturb}}$ in degrees with 95\%
  scene/trajectory-clustered bootstrap intervals; positive values are worse.
  Same-scene query swap is defined for $n=153$ samples from 53 scenes, while
  the other perturbations use all 165 samples.  All probes use training seed 42
  and perturbation seed 42.}
  \label{tab:probe-validity-b2}
  \scriptsize
  \setlength{\tabcolsep}{3.5pt}
  \begin{tabular}{@{}lrrr@{}}
    \toprule
    \multicolumn{4}{@{}l}{\textit{Panel A: query dependence}} \\
    Model & FULL $\downarrow$ & Query mask $\Delta$ & Same-scene swap $\Delta$ \\
    \midrule
    V-JEPA2   & 52.77 & $+15.79\ [ 9.46,21.92]$ & $+11.76\ [ 3.56,17.78]$ \\
    VideoMAE  & 55.77 & $+11.18\ [ 4.88,17.65]$ & $ +9.53\ [ 0.95,16.40]$ \\
    Wan       & 63.50 & $ +4.08\ [-2.58,10.81]$ & $ +2.95\ [-4.88,10.30]$ \\
    CogVideoX & 71.24 & $ -3.98\ [-11.60,3.23]$ & $ +8.29\ [ 1.33,15.20]$ \\
    \midrule
    \multicolumn{4}{@{}l}{\textit{Panel B: history content and content-slot order}} \\
    Model & Current frame $\Delta$ & Remove visible $\Delta$ & Shuffle / reverse $\Delta$ \\
    \midrule
    V-JEPA2   & $+10.26\ [ 3.51,16.80]$ & $+9.22\ [ 3.95,14.59]$ & $+0.83\ [-0.49,2.78]\;/\;+1.03\ [-0.56,3.04]$ \\
    VideoMAE  & $ +2.08\ [-6.10, 9.97]$ & $+5.83\ [-2.09,13.57]$ & $+0.41\ [-0.79,1.67]\;/\;+0.65\ [-0.74,2.13]$ \\
    Wan       & $+10.01\ [ 3.42,16.69]$ & $+6.59\ [ 0.58,12.91]$ & $+0.49\ [-1.66,2.63]\;/\;+1.48\ [-1.08,4.11]$ \\
    CogVideoX & $ -3.41\ [-9.39, 2.39]$ & $+0.54\ [-4.33, 5.53]$ & $-0.07\ [-0.21,0.00]\;/\;-0.06\ [-0.19,0.00]$ \\
    \bottomrule
  \end{tabular}
\end{table*}

\paragraph{Scope and limitations.}
The horizon check is descriptive rather than monotonic: p12 hard-K8 gains for
$h=1/2/4$ are respectively 18.84/18.84/14.49 points for V-JEPA2,
10.51/9.78/15.94 for VideoMAE, 2.90/10.51/3.99 for Wan, and
4.35/6.88/2.17 for CogVideoX.  The full candidate-grid, retrained C1 controls,
and B2 priors remain auxiliary because their execution or fitting provenance is
incomplete; in particular, the prior analysis does not exhaust category- or
layout-conditioned alternatives.  All analyses use a single probe-training
seed; candidate seeds are sampling seeds, not additional training seeds.  The
results are development-validation analyses without an independent frozen
confirmation.  The nominal master test is not independent of the historical
development universe, but this does not demonstrate supervised sample leakage
between the actual legacy C train and validation partitions.  Finally, p12 and
p16 are not a paired history-length experiment, model architecture and training
data are not controlled for objective-level attribution, and B2/P16
perturbations were not verified for V-JEPA2, VideoMAE, or Wan.  Real alternative
futures and order-critical paired trajectories are outside the present data.
